% Options for packages loaded elsewhere
\PassOptionsToPackage{unicode}{hyperref}
\PassOptionsToPackage{hyphens}{url}
%
\documentclass[
]{article}
\usepackage{amsmath,amssymb}
\usepackage{iftex}
\ifPDFTeX
  \usepackage[T1]{fontenc}
  \usepackage[utf8]{inputenc}
  \usepackage{textcomp} % provide euro and other symbols
\else % if luatex or xetex
  \usepackage{unicode-math} % this also loads fontspec
  \defaultfontfeatures{Scale=MatchLowercase}
  \defaultfontfeatures[\rmfamily]{Ligatures=TeX,Scale=1}
\fi
\usepackage{lmodern}
\ifPDFTeX\else
  % xetex/luatex font selection
\fi
% Use upquote if available, for straight quotes in verbatim environments
\IfFileExists{upquote.sty}{\usepackage{upquote}}{}
\IfFileExists{microtype.sty}{% use microtype if available
  \usepackage[]{microtype}
  \UseMicrotypeSet[protrusion]{basicmath} % disable protrusion for tt fonts
}{}
\makeatletter
\@ifundefined{KOMAClassName}{% if non-KOMA class
  \IfFileExists{parskip.sty}{%
    \usepackage{parskip}
  }{% else
    \setlength{\parindent}{0pt}
    \setlength{\parskip}{6pt plus 2pt minus 1pt}}
}{% if KOMA class
  \KOMAoptions{parskip=half}}
\makeatother
\usepackage{xcolor}
\usepackage{longtable,booktabs,array}
\usepackage{calc} % for calculating minipage widths
% Correct order of tables after \paragraph or \subparagraph
\usepackage{etoolbox}
\makeatletter
\patchcmd\longtable{\par}{\if@noskipsec\mbox{}\fi\par}{}{}
\makeatother
% Allow footnotes in longtable head/foot
\IfFileExists{footnotehyper.sty}{\usepackage{footnotehyper}}{\usepackage{footnote}}
\makesavenoteenv{longtable}
\setlength{\emergencystretch}{3em} % prevent overfull lines
\providecommand{\tightlist}{%
  \setlength{\itemsep}{0pt}\setlength{\parskip}{0pt}}
\setcounter{secnumdepth}{-\maxdimen} % remove section numbering
\ifLuaTeX
  \usepackage{selnolig}  % disable illegal ligatures
\fi
\IfFileExists{bookmark.sty}{\usepackage{bookmark}}{\usepackage{hyperref}}
\IfFileExists{xurl.sty}{\usepackage{xurl}}{} % add URL line breaks if available
\urlstyle{same}
\hypersetup{
  hidelinks,
  pdfcreator={LaTeX via pandoc}}

\author{}
\date{}

\begin{document}

\hypertarget{ayudantuxeda-6-soluciuxf3n---planificaciuxf3n-agregada-y-mrp}{%
\section{Ayudantía 6 Solución - Planificación agregada y
MRP}\label{ayudantuxeda-6-soluciuxf3n---planificaciuxf3n-agregada-y-mrp}}

\textbf{Pontificia Universidad Católica de Chile}\\
\textbf{Escuela de Ingeniería}\\
\textbf{Departamento de Ingeniería Industrial y de Sistemas}\\
\textbf{ICS3213: Gestión de Operaciones}

\begin{center}\rule{0.5\linewidth}{0.5pt}\end{center}

\hypertarget{problema-1-tarea-ii-2024---soluciuxf3n}{%
\subsection{Problema 1 (Tarea II 2024) -
Solución}\label{problema-1-tarea-ii-2024---soluciuxf3n}}

\hypertarget{a-modelo-de-programaciuxf3n-lineal-entera}{%
\subsubsection{a) Modelo de Programación Lineal
Entera}\label{a-modelo-de-programaciuxf3n-lineal-entera}}

\hypertarget{conjuntos}{%
\paragraph{Conjuntos}\label{conjuntos}}

\begin{itemize}
\tightlist
\item
  \(T = \{1, 2, \dots, 12\}\): Conjunto de períodos (meses de la
  planificación anual).
\end{itemize}

\hypertarget{paruxe1metros}{%
\paragraph{Parámetros}\label{paruxe1metros}}

\begin{itemize}
\tightlist
\item
  \(d_t\): Demanda de cajas de suministro en el período \(t \in T\).
\item
  \(l_{so}\): Límite de demanda insatisfecha (faltantes) permitido por
  período.
\item
  \(c_{so}\): Costo de penalización por cada caja de suministros no
  entregada a tiempo.
\item
  \(c_p\): Costo base de producir una caja de suministros de manera
  propia.
\item
  \(c_{su}\): Costo de subcontratar la producción de una caja de
  suministros.
\item
  \(l_{su}\): Límite de unidades que se pueden subcontratar por período.
\item
  \(i_0\): Inventario inicial de cajas de suministro.
\item
  \(c_i\): Costo de mantener una caja de suministro en el inventario por
  período.
\item
  \(w_0\): Dotación inicial (y final requerida) de trabajadores.
\item
  \(c_c\): Costo de contratar un trabajador.
\item
  \(c_d\): Costo de despedir un trabajador.
\item
  \(h\): Duración de la jornada diaria de trabajo (horas/turno).
\item
  \(n_t\): Número de días de trabajo en el mes \(t\).
\item
  \(c_{sn}\): Salario por hora normal de trabajo.
\item
  \(l_{he}\): Límite de horas extra permitidas por trabajador por
  período.
\item
  \(c_{se}\): Salario por hora extra de trabajo.
\item
  \(\alpha\): Factor de productividad de un trabajador novato
  (\(0 \le \alpha \le 1\)).
\item
  \(p_{hn}\): Productividad de un trabajador experimentado en hora
  normal (cajas/HH).
\item
  \(p_{he}\): Productividad de un trabajador experimentado en hora extra
  (cajas/HH).
\end{itemize}

\hypertarget{variables-de-decisiuxf3n}{%
\paragraph{Variables de Decisión}\label{variables-de-decisiuxf3n}}

\begin{itemize}
\tightlist
\item
  \(W_t\): Dotación total de trabajadores en el período \(t \in T\)
  (\(W_t \in \mathbb{Z}^+\)).
\item
  \(C_t\): Cantidad de trabajadores contratados al inicio del período
  \(t \in T\) (\(C_t \in \mathbb{Z}^+\)).
\item
  \(D_t\): Cantidad de trabajadores despedidos al inicio del período
  \(t \in T\) (\(D_t \in \mathbb{Z}^+\)).
\item
  \(P_t\): Cantidad de cajas de suministros producidas de forma interna
  en el período \(t \in T\) (\(P_t \ge 0\)).
\item
  \(SU_t\): Cantidad de cajas de suministros producidas por
  subcontratación en el período \(t \in T\) (\(SU_t \ge 0\)).
\item
  \(PD_t\): Cantidad de cajas de suministros despachadas (entregadas) en
  el período \(t \in T\) (\(PD_t \ge 0\)).
\item
  \(SO_t\): Cantidad de cajas de suministros de demanda insatisfecha
  (faltantes) al final del período \(t \in T\) (\(SO_t \ge 0\)).
\item
  \(HE_t\): Horas extras totales trabajadas en la planta en el período
  \(t \in T\) (\(HE_t \ge 0\)).
\item
  \(I_t\): Cantidad de cajas de suministros en el inventario al final
  del período \(t \in T\) (\(I_t \ge 0\)).
\end{itemize}

\hypertarget{funciuxf3n-objetivo}{%
\paragraph{Función Objetivo}\label{funciuxf3n-objetivo}}

Minimizar los costos totales, los cuales incluyen: subcontratación,
penalización por faltantes, contratación, despidos, salarios normales,
horas extra, almacenamiento e inventario, y costo de producción propia
(con descuento del \(30\%\) durante los primeros 3 meses):

\[\begin{aligned}
\min \sum_{t=1}^{12} \Big( &c_{su} SU_t + c_{so} SO_t + c_c C_t + c_d D_t + c_{sn} \cdot h \cdot n_t \cdot W_t + c_{se} HE_t + c_i I_t + \text{costo\_prod}_t \Big)
\end{aligned}\]

Donde el costo de producción propia es:
\[\text{costo\_prod}_t = \begin{cases} 
0.7 c_p P_t & \text{si } t \in \{1, 2, 3\} \\
c_p P_t & \text{si } t \in \{4, \dots, 12\}
\end{cases}\]

\hypertarget{restricciones}{%
\paragraph{Restricciones}\label{restricciones}}

\begin{enumerate}
\def\labelenumi{\arabic{enumi}.}
\item
  \textbf{Flujo de trabajadores:} \[W_1 = w_0 + C_1 - D_1\]
  \[W_t = W_{t-1} + C_t - D_t \quad \forall t \in \{2, \dots, 12\}\]
  \[W_{12} = w_0\]
\item
  \textbf{Flujo de inventarios:} \[I_1 = i_0 + P_1 + SU_1 - PD_1\]
  \[I_t = I_{t-1} + P_t + SU_t - PD_t \quad \forall t \in \{2, \dots, 12\}\]
\item
  \textbf{Capacidad de producción propia:} Considerando que los
  trabajadores contratados en el período \(t\) son novatos con
  productividad \(\alpha \cdot p_{hn}\) en horario normal, y que los
  experimentados (\(W_t - C_t\)) rinden a \(p_{hn}\):
  \[P_t \le \Big[ (W_t - C_t) + \alpha C_t \Big] \cdot h \cdot n_t \cdot p_{hn} + HE_t \cdot p_{he} \quad \forall t \in T\]
\item
  \textbf{Despacho y demanda insatisfecha:} La cantidad despachada es la
  demanda menos la demanda insatisfecha acumulada del período actual,
  considerando que lo no entregado se arrastra:
  \[PD_t = d_t - SO_t \quad \forall t \in T\]
\item
  \textbf{Límite de horas extras:}
  \[HE_t \le l_{he} \cdot W_t \quad \forall t \in T\]
\item
  \textbf{Límite de producción subcontratada:}
  \[SU_t \le l_{su} \quad \forall t \in T\]
\item
  \textbf{Límite de demanda insatisfecha:}
  \[SO_t \le l_{so} \quad \forall t \in T\]
\item
  \textbf{Naturaleza de variables:}
  \[W_t, C_t, D_t \in \mathbb{Z}^+ \quad \forall t \in T\]
  \[P_t, SU_t, PD_t, SO_t, HE_t, I_t \ge 0 \quad \forall t \in T\]
\end{enumerate}

\begin{center}\rule{0.5\linewidth}{0.5pt}\end{center}

\hypertarget{b-restricciuxf3n-de-capacidad-de-almacenamiento-limitada}{%
\subsubsection{b) Restricción de Capacidad de Almacenamiento
Limitada}\label{b-restricciuxf3n-de-capacidad-de-almacenamiento-limitada}}

Se agrega el parámetro \(I_{max} = 50\) (inventario máximo permitido por
período) y las siguientes restricciones:
\[I_t \le 50 \quad \forall t \in T\]

\begin{center}\rule{0.5\linewidth}{0.5pt}\end{center}

\hypertarget{c-sin-subcontrataciuxf3n}{%
\subsubsection{c) Sin Subcontratación}\label{c-sin-subcontrataciuxf3n}}

Se elimina la posibilidad de subcontratar. Esto implica: 1. Eliminar la
variable \(SU_t\) del modelo (o fijar
\(SU_t = 0 \quad \forall t \in T\)). 2. Eliminar el término
\(c_{su} SU_t\) de la función objetivo. 3. Eliminar la restricción de
límite de producción subcontratada: \(SU_t \le l_{su}\).

\begin{center}\rule{0.5\linewidth}{0.5pt}\end{center}

\hypertarget{problema-2-i2-2024---soluciuxf3n}{%
\subsection{Problema 2 (I2 2024) -
Solución}\label{problema-2-i2-2024---soluciuxf3n}}

\hypertarget{a-lista-de-materiales-bom}{%
\subsubsection{a) Lista de Materiales
(BOM)}\label{a-lista-de-materiales-bom}}

\begin{verbatim}
Producto Final (1 envase)
├── Envase (1 unidad)
├── Compuesto A (300 ml = 0.3 L)
└── Compuesto B (200 ml = 0.2 L)
    ├── Compuesto C (100 ml = 0.1 L)
    └── Compuesto D (100 ml = 0.1 L)
\end{verbatim}

\emph{Nota:} Dado que la receta está en ml y los inventarios se miden en
litros para los compuestos químicos, utilizaremos la conversión
\(1\text{ Litro} = 1000\text{ ml}\).

\begin{itemize}
\tightlist
\item
  1 unidad de Producto Final requiere:

  \begin{itemize}
  \tightlist
  \item
    1 pomo de Envase
  \item
    \(0.3\text{ L}\) de Compuesto A
  \item
    \(0.2\text{ L}\) de Compuesto B (que a su vez requiere
    \(0.1\text{ L}\) de C y \(0.1\text{ L}\) de D).
  \end{itemize}
\end{itemize}

\begin{center}\rule{0.5\linewidth}{0.5pt}\end{center}

\hypertarget{b-tablas-de-mrp}{%
\subsubsection{b) Tablas de MRP}\label{b-tablas-de-mrp}}

\hypertarget{producto-final-lead-time-1-semana-lote-a-lote}{%
\paragraph{1. Producto Final (Lead Time = 1 semana, Lote a
Lote)}\label{producto-final-lead-time-1-semana-lote-a-lote}}

\begin{longtable}[]{@{}lccccccc@{}}
\toprule\noalign{}
Semana & 0 & 1 & 2 & 3 & 4 & 5 & 6 \\
\midrule\noalign{}
\endhead
\bottomrule\noalign{}
\endlastfoot
\textbf{Requerimiento Bruto} & & 300 & 200 & 300 & 400 & 300 & 200 \\
\textbf{Inventario Disponible} & 500 & 200 & 0 & 0 & 0 & 0 & 0 \\
\textbf{Requerimiento Neto} & & 0 & 0 & 100 & 400 & 300 & 200 \\
\textbf{Recepción Planeada} & & 0 & 0 & 100 & 400 & 300 & 200 \\
\textbf{Lanzamiento Planeado} & & 0 & 100 & 400 & 300 & 200 & 0 \\
\end{longtable}

\hypertarget{envase-lead-time-3-semanas-lote-a-lote}{%
\paragraph{2. Envase (Lead Time = 3 semanas, Lote a
Lote)}\label{envase-lead-time-3-semanas-lote-a-lote}}

\emph{Requerimiento bruto = Lanzamiento planeado de Producto Final (1
unidad por unidad).}

\begin{longtable}[]{@{}lccccccc@{}}
\toprule\noalign{}
Semana & 0 & 1 & 2 & 3 & 4 & 5 & 6 \\
\midrule\noalign{}
\endhead
\bottomrule\noalign{}
\endlastfoot
\textbf{Requerimiento Bruto} & & 0 & 100 & 400 & 300 & 200 & 0 \\
\textbf{Recepciones Programadas} & & 0 & 400 & 0 & 0 & 0 & 0 \\
\textbf{Inventario Disponible} & 300 & 300 & 600 & 200 & 0 & 0 & 0 \\
\textbf{Requerimiento Neto} & & 0 & 0 & 0 & 100 & 200 & 0 \\
\textbf{Recepción Planeada} & & 0 & 0 & 0 & 100 & 200 & 0 \\
\textbf{Lanzamiento Planeado} & 100 & 200 & 0 & 0 & 0 & 0 & 0 \\
\end{longtable}

\emph{Nota:} El lanzamiento de 100 en la semana 0 y 200 en la semana 1
debe programarse en el pasado.

\hypertarget{compuesto-a-lead-time-1-semana-lote-a-lote}{%
\paragraph{3. Compuesto A (Lead Time = 1 semana, Lote a
Lote)}\label{compuesto-a-lead-time-1-semana-lote-a-lote}}

\emph{Requerimiento bruto = Lanzamiento planeado de Producto Final
\(\times 0.3\text{ L}\).}

\begin{longtable}[]{@{}lccccccc@{}}
\toprule\noalign{}
Semana & 0 & 1 & 2 & 3 & 4 & 5 & 6 \\
\midrule\noalign{}
\endhead
\bottomrule\noalign{}
\endlastfoot
\textbf{Requerimiento Bruto} & & 0 & 30 & 120 & 90 & 60 & 0 \\
\textbf{Inventario Disponible} & 0 & 0 & 0 & 0 & 0 & 0 & 0 \\
\textbf{Requerimiento Neto} & & 0 & 30 & 120 & 90 & 60 & 0 \\
\textbf{Recepción Planeada} & & 0 & 30 & 120 & 90 & 60 & 0 \\
\textbf{Lanzamiento Planeado} & 0 & 30 & 120 & 90 & 60 & 0 & 0 \\
\end{longtable}

\hypertarget{compuesto-b-lead-time-1-semana-lote-a-lote}{%
\paragraph{4. Compuesto B (Lead Time = 1 semana, Lote a
Lote)}\label{compuesto-b-lead-time-1-semana-lote-a-lote}}

\emph{Requerimiento bruto = Lanzamiento planeado de Producto Final
\(\times 0.2\text{ L}\).}

\begin{longtable}[]{@{}lccccccc@{}}
\toprule\noalign{}
Semana & 0 & 1 & 2 & 3 & 4 & 5 & 6 \\
\midrule\noalign{}
\endhead
\bottomrule\noalign{}
\endlastfoot
\textbf{Requerimiento Bruto} & & 0 & 20 & 80 & 60 & 40 & 0 \\
\textbf{Inventario Disponible} & 50 & 50 & 30 & 0 & 0 & 0 & 0 \\
\textbf{Requerimiento Neto} & & 0 & 0 & 50 & 60 & 40 & 0 \\
\textbf{Recepción Planeada} & & 0 & 0 & 50 & 60 & 40 & 0 \\
\textbf{Lanzamiento Planeado} & 0 & 0 & 50 & 60 & 40 & 0 & 0 \\
\end{longtable}

\hypertarget{compuesto-c-lead-time-2-semanas-lote-a-lote}{%
\paragraph{5. Compuesto C (Lead Time = 2 semanas, Lote a
Lote)}\label{compuesto-c-lead-time-2-semanas-lote-a-lote}}

\emph{Requerimiento bruto = Lanzamiento planeado de Compuesto B
\(\times (100\text{ ml} / 200\text{ ml}) = 0.5\text{ L}\) por litro de
B.} \emph{Es decir, \(50\%\) del compuesto B es compuesto C.} * Semana
2: Req. Bruto = \(50 \times 0.5 = 25\text{ L}\). * Semana 3: Req. Bruto
= \(60 \times 0.5 = 30\text{ L}\). * Semana 4: Req. Bruto =
\(40 \times 0.5 = 20\text{ L}\).

\begin{longtable}[]{@{}lccccccc@{}}
\toprule\noalign{}
Semana & 0 & 1 & 2 & 3 & 4 & 5 & 6 \\
\midrule\noalign{}
\endhead
\bottomrule\noalign{}
\endlastfoot
\textbf{Requerimiento Bruto} & & 0 & 25 & 30 & 20 & 0 & 0 \\
\textbf{Inventario Disponible} & 50 & 50 & 25 & 0 & 0 & 0 & 0 \\
\textbf{Requerimiento Neto} & & 0 & 0 & 5 & 20 & 0 & 0 \\
\textbf{Recepción Planeada} & & 0 & 0 & 5 & 20 & 0 & 0 \\
\textbf{Lanzamiento Planeado} & 0 & 5 & 20 & 0 & 0 & 0 & 0 \\
\end{longtable}

\hypertarget{compuesto-d-lead-time-1-semana-lote-a-lote}{%
\paragraph{6. Compuesto D (Lead Time = 1 semana, Lote a
Lote)}\label{compuesto-d-lead-time-1-semana-lote-a-lote}}

\emph{Requerimiento bruto = Lanzamiento planeado de Compuesto B
\(\times 0.5\text{ L}\) por litro de B.}

\begin{longtable}[]{@{}lccccccc@{}}
\toprule\noalign{}
Semana & 0 & 1 & 2 & 3 & 4 & 5 & 6 \\
\midrule\noalign{}
\endhead
\bottomrule\noalign{}
\endlastfoot
\textbf{Requerimiento Bruto} & & 0 & 25 & 30 & 20 & 0 & 0 \\
\textbf{Inventario Disponible} & 0 & 0 & 0 & 0 & 0 & 0 & 0 \\
\textbf{Requerimiento Neto} & & 0 & 25 & 30 & 20 & 0 & 0 \\
\textbf{Recepción Planeada} & & 0 & 25 & 30 & 20 & 0 & 0 \\
\textbf{Lanzamiento Planeado} & 0 & 25 & 30 & 20 & 0 & 0 & 0 \\
\end{longtable}

\begin{center}\rule{0.5\linewidth}{0.5pt}\end{center}

\hypertarget{c-algoritmo-de-loteo-para-compuesto-c}{%
\subsubsection{c) Algoritmo de Loteo para Compuesto
C}\label{c-algoritmo-de-loteo-para-compuesto-c}}

Tenemos los siguientes requerimientos netos para el Compuesto C: *
Semana 3: \(5\text{ Litros}\) * Semana 4: \(20\text{ Litros}\)

\textbf{Parámetros:} * Costo de ordenar (Setup) = \(S = 10\) * Costo de
mantener inventario = \(H = 1\) por litro por semana.

Evaluamos las dos políticas posibles:

\hypertarget{opciuxf3n-1-lote-a-lote-l-times-l---no-agrupar}{%
\paragraph{\texorpdfstring{Opción 1: Lote a Lote (\(L \times L\)) - No
agrupar}{Opción 1: Lote a Lote (L \textbackslash times L) - No agrupar}}\label{opciuxf3n-1-lote-a-lote-l-times-l---no-agrupar}}

\begin{itemize}
\tightlist
\item
  \textbf{Semana 3:} Ordenamos \(5\text{ L}\). Costo de ordenar =
  \(10\). Costo de inventario = \(0\).
\item
  \textbf{Semana 4:} Ordenamos \(20\text{ L}\). Costo de ordenar =
  \(10\). Costo de inventario = \(0\).
\item
  \textbf{Costo Total = 20.}
\end{itemize}

\hypertarget{opciuxf3n-2-agrupar-en-la-semana-3-la-demanda-de-la-semana-3-y-4-ordenar-25text-l-en-semana-3}{%
\paragraph{\texorpdfstring{Opción 2: Agrupar en la Semana 3 la demanda
de la semana 3 y 4 (Ordenar \(25\text{ L}\) en Semana
3)}{Opción 2: Agrupar en la Semana 3 la demanda de la semana 3 y 4 (Ordenar 25\textbackslash text\{ L\} en Semana 3)}}\label{opciuxf3n-2-agrupar-en-la-semana-3-la-demanda-de-la-semana-3-y-4-ordenar-25text-l-en-semana-3}}

\begin{itemize}
\tightlist
\item
  \textbf{Semana 3:} Recibimos \(25\text{ L}\). Costo de ordenar =
  \(10\).
\item
  Mantenemos \(20\text{ L}\) en inventario de la semana 3 a la 4 (1
  semana).
\item
  Costo de mantener inventario =
  \(20\text{ L} \times 1\text{ semana} \times \$1 = 20\).
\item
  \textbf{Costo Total = 10 (orden) + 20 (inventario) = 30.}
\end{itemize}

\hypertarget{conclusiuxf3n}{%
\paragraph{Conclusión:}\label{conclusiuxf3n}}

La mejor opción es la \textbf{Opción 1: Lote a Lote} con un costo de
\textbf{\(20\)}, comparado con los \(30\) de agrupar. No conviene
agrupar lotes debido a que el costo de mantener inventario (\(20\))
supera el ahorro de evitar un setup (\(10\)).

\begin{center}\rule{0.5\linewidth}{0.5pt}\end{center}

\hypertarget{problema-3-pert---soluciuxf3n}{%
\subsection{Problema 3 (PERT) -
Solución}\label{problema-3-pert---soluciuxf3n}}

\hypertarget{tiempos-tempranos-tarduxedos-y-ruta-cruxedtica}{%
\subsubsection{1. Tiempos tempranos, tardíos y Ruta
Crítica}\label{tiempos-tempranos-tarduxedos-y-ruta-cruxedtica}}

Para cada actividad se calculan el tiempo esperado y la varianza: *
\(t_e = \frac{a + 4m + b}{6}\) *
\(\sigma^2 = \left(\frac{b - a}{6}\right)^2\)

\begin{longtable}[]{@{}ccccccc@{}}
\toprule\noalign{}
Actividad & Predecesores & \(a\) & \(m\) & \(b\) & \(t_e\) &
\(\sigma^2\) \\
\midrule\noalign{}
\endhead
\bottomrule\noalign{}
\endlastfoot
\textbf{A} & --- & 4 & 5 & 6 & 5 & 0.09 \\
\textbf{B} & --- & 2 & 3 & 4 & 3 & 0.04 \\
\textbf{C} & A, B & 3 & 4 & 5 & 4 & 0.09 \\
\textbf{D} & C & 2 & 3 & 4 & 3 & 0.09 \\
\textbf{E} & C & 3 & 5 & 7 & 5 & 0.44 \\
\textbf{F} & D, E & 2 & 4 & 6 & 4 & 0.44 \\
\end{longtable}

\hypertarget{cuxe1lculo-de-calendario}{%
\paragraph{Cálculo de Calendario:}\label{cuxe1lculo-de-calendario}}

\begin{longtable}[]{@{}ccccccc@{}}
\toprule\noalign{}
Actividad & \(t_e\) & ES & EF & LS & LF & Holgura (\(LF - EF\)) \\
\midrule\noalign{}
\endhead
\bottomrule\noalign{}
\endlastfoot
\textbf{A} & 5 & 0 & 5 & 0 & 5 & 0 (Crítica) \\
\textbf{B} & 3 & 0 & 3 & 2 & 5 & 2 \\
\textbf{C} & 4 & 5 & 9 & 5 & 9 & 0 (Crítica) \\
\textbf{D} & 3 & 9 & 12 & 11 & 14 & 2 \\
\textbf{E} & 5 & 9 & 14 & 9 & 14 & 0 (Crítica) \\
\textbf{F} & 4 & 14 & 18 & 14 & 18 & 0 (Crítica) \\
\end{longtable}

\begin{itemize}
\tightlist
\item
  \textbf{Ruta Crítica:} \(A \rightarrow C \rightarrow E \rightarrow F\)
\item
  \textbf{Duración Esperada (\(T_e\)):} \(5 + 4 + 5 + 4 = 18\) semanas.
\item
  \textbf{Desviación Estándar de la Ruta Crítica (\(\sigma_T\)):}
  \[\sigma_T = \sqrt{\sigma_A^2 + \sigma_C^2 + \sigma_E^2 + \sigma_F^2} = \sqrt{0.09 + 0.09 + 0.44 + 0.44} = \sqrt{1.06} \approx 1.03 \text{ semanas}\]
  \emph{(Nota: Si usamos la desviación estándar directa del cuadro
  original \(\sigma_i\):
  \(\sigma_T = \sqrt{0.9^2 + 1.0^2 + 1.1^2 + 0.8^2} = \sqrt{0.81 + 1.0 + 1.21 + 0.64} = \sqrt{3.66} \approx 1.91\)
  semanas).}
\end{itemize}

\begin{center}\rule{0.5\linewidth}{0.5pt}\end{center}

\hypertarget{duraciuxf3n-doble-de-probabilidad-de-excederse-que-de-cumplir}{%
\subsubsection{2. Duración doble de probabilidad de excederse que de
cumplir}\label{duraciuxf3n-doble-de-probabilidad-de-excederse-que-de-cumplir}}

Sea \(T_0\) el plazo buscado. Queremos que:
\[P(T > T_0) = 2 P(T \le T_0)\] Como \(P(T > T_0) + P(T \le T_0) = 1\),
sustituyendo obtenemos:
\[3 P(T \le T_0) = 1 \implies P(T \le T_0) = 0.3333\]

Buscando en la tabla de la normal estándar para una probabilidad
acumulada de \(0.3333\), obtenemos un valor de \(Z \approx -0.43\).
\[\frac{T_0 - 18}{1.91} = -0.43 \implies T_0 = 18 - 0.43 \times 1.91 \approx 17.18 \text{ semanas}\]

\begin{center}\rule{0.5\linewidth}{0.5pt}\end{center}

\hypertarget{contrato-con-bono-y-penalizaciuxf3n}{%
\subsubsection{3. Contrato con Bono y
Penalización}\label{contrato-con-bono-y-penalizaciuxf3n}}

\begin{itemize}
\tightlist
\item
  \textbf{Bono:} \(\$200\) si se termina en \(\le 15\) semanas.
\item
  \textbf{Penalización:} \(-\$80\) si se termina en \(\ge 20\) semanas.
\end{itemize}

\hypertarget{probabilidad-de-ganar-el-bono-t-le-15}{%
\paragraph{\texorpdfstring{Probabilidad de ganar el Bono
(\(T \le 15\)):}{Probabilidad de ganar el Bono (T \textbackslash le 15):}}\label{probabilidad-de-ganar-el-bono-t-le-15}}

\[Z_{bono} = \frac{15 - 18}{1.91} = -1.57 \implies P(T \le 15) = P(Z \le -1.57) \approx 0.0582 \text{ (5.8\%)}\]

\hypertarget{probabilidad-de-pagar-penalizaciuxf3n-t-ge-20}{%
\paragraph{\texorpdfstring{Probabilidad de pagar Penalización
(\(T \ge 20\)):}{Probabilidad de pagar Penalización (T \textbackslash ge 20):}}\label{probabilidad-de-pagar-penalizaciuxf3n-t-ge-20}}

\[Z_{pen} = \frac{20 - 18}{1.91} = 1.05 \implies P(T \ge 20) = P(Z \ge 1.05) = 1 - 0.8531 = 0.1469 \text{ (14.7\%)}\]

\hypertarget{valor-esperado-ev}{%
\paragraph{Valor Esperado (EV):}\label{valor-esperado-ev}}

\[EV = 200 \times P(T \le 15) - 80 \times P(T \ge 20) = 200 \times 0.0582 - 80 \times 0.1469 = 11.64 - 11.75 = -0.11\]

Como el \textbf{Valor Esperado es negativo (-\$0.11)}, se debería
\textbf{rechazar} el contrato bajo estas condiciones.

\hypertarget{plazo-muxe1ximo-del-bono-para-aceptar}{%
\paragraph{Plazo máximo del bono para
aceptar:}\label{plazo-muxe1ximo-del-bono-para-aceptar}}

Para que \(EV > 0\):
\[200 \times P(T \le X) - 80 \times 0.1469 = 0 \implies P(T \le X) = \frac{11.75}{200} = 0.0588\]
Esto da un \(Z \approx -1.56\).
\[\frac{X - 18}{1.91} = -1.56 \implies X \approx 15.02 \text{ semanas}\]
El plazo máximo para ofrecer el bono debe ser de al menos \(15.02\)
semanas para que el valor esperado sea positivo.

\begin{center}\rule{0.5\linewidth}{0.5pt}\end{center}

\hypertarget{reducciuxf3n-de-3-semanas-con-costo-muxednimo}{%
\subsubsection{4. Reducción de 3 semanas con costo
mínimo}\label{reducciuxf3n-de-3-semanas-con-costo-muxednimo}}

Queremos reducir el proyecto en 3 semanas. Las actividades en la ruta
crítica que podemos acelerar a su tiempo optimista son:

\begin{itemize}
\tightlist
\item
  \textbf{Actividad A:} De \(5 \rightarrow 4\) semanas. Ahorro = \(1\)
  semana. Costo adicional = \(\$10\) (el doble del costo normal). Costo
  por semana de aceleración = \(\$10\).
\item
  \textbf{Actividad C:} De \(4 \rightarrow 3\) semanas. Ahorro = \(1\)
  semana. Costo adicional = \(\$7\). Costo por semana de aceleración =
  \(\$7\).
\item
  \textbf{Actividad E:} De \(5 \rightarrow 3\) semanas. Ahorro = \(2\)
  semanas. Costo adicional = \(\$12\). Costo por semana de aceleración =
  \(\$6\).
\item
  \textbf{Actividad F:} De \(4 \rightarrow 2\) semanas. Ahorro = \(2\)
  semanas. Costo adicional = \(\$10\). Costo por semana de aceleración =
  \(\$5\).
\end{itemize}

\hypertarget{estrategia-de-aceleraciuxf3n-al-costo-muxednimo}{%
\paragraph{Estrategia de Aceleración al Costo
Mínimo:}\label{estrategia-de-aceleraciuxf3n-al-costo-muxednimo}}

\begin{enumerate}
\def\labelenumi{\arabic{enumi}.}
\tightlist
\item
  Aceleramos \textbf{F} por 2 semanas.

  \begin{itemize}
  \tightlist
  \item
    Costo de aceleración = \(\$10\).
  \item
    Duración acumulada reducida = 2 semanas.
  \end{itemize}
\item
  Para la tercera semana de reducción, la opción más barata disponible
  en la ruta crítica es acelerar \textbf{E} (costo de \(\$6\) por
  semana) o \textbf{C} (costo de \(\$7\) por semana).

  \begin{itemize}
  \tightlist
  \item
    Dado que el enunciado indica que no se permite aceleración parcial,
    debemos acelerar \textbf{C} por 1 semana (costo \(\$7\)).
  \item
    \emph{Nota:} Si aceleráramos E, tendríamos que reducir las 2 semanas
    completas con un costo de \(\$12\), lo cual es más caro que acelerar
    C por \(\$7\).
  \end{itemize}
\item
  \textbf{Costo Mínimo Adicional:}
  \(\$10 \text{ (por F)} + \$7 \text{ (por C)} = \$17\).
\item
  \textbf{Costo Total del Proyecto:} Costo normal total
  \((\$26) + \$17 = \$43\).
\end{enumerate}

\end{document}
